\chapter{MÃ HÓA VÀ KẾ HOẠCH PHÂN TÍCH}
\section*{Quy tắc mã hóa}
Các biến Likert nhận giá trị 1--5 theo thứ tự tăng mức đồng ý. Giá trị thiếu để trống, không mã hóa bằng 0. Biến giới tính, địa bàn, khu vực và học lực dùng mã có nhãn; lưu riêng từ điển dữ liệu. Điểm nhân tố được tính bằng trung bình các biến giữ lại sau kiểm định, với điều kiện người trả lời hoàn thành ít nhất 75\% số mục trong nhân tố.

\section*{Chuỗi phân tích tái lập}
\begin{enumerate}
\item Khóa tệp dữ liệu gốc ở chế độ chỉ đọc; tạo tệp làm sạch bằng mã lệnh.
\item Báo cáo số phiếu loại và lý do; kiểm tra thiếu, ngoại lệ, mẫu trả lời bất thường.
\item Thống kê đặc điểm mẫu và từng biến quan sát.
\item Kiểm định Alpha/Omega; EFA bằng Principal Axis Factoring và Promax; dùng parallel analysis nếu phần mềm cho phép.
\item Tính điểm nhân tố, kiểm tra tương quan và giả định hồi quy.
\item Ước lượng mô hình chính và mô hình độ bền với sai số chuẩn HC3.
\item So sánh nhóm kèm kích thước ảnh hưởng và khoảng tin cậy.
\item Xuất bảng từ mã lệnh để hạn chế sao chép thủ công; ẩn mọi thông tin nhận diện.
\end{enumerate}

\section*{Danh mục hồ sơ cần lưu}
Đề cương đã duyệt; văn bản cho phép khảo sát; biểu mẫu đồng thuận; bảng hỏi cuối; dữ liệu gốc ẩn danh; nhật ký làm sạch; mã phân tích; đầu ra phần mềm; bảng/hình dùng trong luận văn. Thời hạn lưu và quyền truy cập tuân theo quy định của Trường Đại học Cần Thơ và đơn vị đào tạo.

